%paper also and it is very nice !
\documentclass[a4paper,11pt]{article}

\usepackage[a4paper,margin=1in]{geometry}
%\usepackage[a4paper,margin=1in,landscape]{geometry}
\usepackage{graphicx}
%\usepackage{float}% If comment this, figure moves
\usepackage{mwe}    % loads »blindtext« and »graphicx«
%\usepackage{subfig} % error
\usepackage{subcaption} % sub figure placement
\usepackage[english]{babel}
\usepackage{ucs}               % require separate debian package 
\usepackage[utf8x]{inputenc}   % require ucs. Allows both Italian and English letters.
\usepackage{textcomp}          % allow euro symbol
\usepackage{fancyhdr}
\usepackage{amssymb,amsmath}
\usepackage{threeparttable}
\usepackage{color}
\usepackage{url}
\usepackage{natbib}  % substituted hardvard with natbib as the former package doesn't like the hyperref one
\usepackage{rotating} %to have text vertical alignment in some table headers
\usepackage{url}
\usepackage{moreverb}
\usepackage{tabularx} % allow line breaks in tables
\usepackage{html,htmllist} % don't think will be needed
\usepackage{hyperref}
\usepackage{mhchem} % for \ce{CO2}

\setlength{\pdfpagewidth}{\paperwidth}
\setlength{\pdfpageheight}{\paperheight}

\bibliographystyle{agsmurl}  % bibliographic style


%opening
\title{FFSM++ complete output, Oracle project}

\author{Antonello Lobianco, Sylvain Caurla, Philippe Delacote, Jean-Daniel Bontemps, Pierre Mérian, Anna Lungarska, Ahmed Barkaoui}

\date{\today}

\bibliographystyle{agsmurl}

\begin{document}
\maketitle

\begin{abstract}
This book contains the complete results of FFSM++ (within the project Oracle) in terms of tables and charts, hierarchically organised by region and theme.

Tables present a comparison (either at 2100 or as average of multiple years) between a reference scenario and a comparing one. If the latter one is a random scenario run multiple times some basic statistic is given.

On the other side charts provide a quick overview of the dynamic of key variables over time.

This document is generated automatically trough a script directly from the model's output. Please refer to the article for model description, scenario definition and output interpretation.

\end{abstract}

\thispagestyle{empty}
\newpage
\tableofcontents
\thispagestyle{empty}
\newpage
\part[Tables]{Output tables}
\input{tables/00_master_tables}
\clearpage
\newpage
\part[Charts]{Output charts}
\input{charts/00_master_charts}
\clearpage
\newpage
%\part[Maps]{Output maps}
%\input{maps/bitmaps/00_master_maps}

%\newpage
%\bibliography{biblio}

\end{document}
\endinput

